% GENERATED FILE. Edit canonical lessons, then run npm run generate:books.
% canonical-source-hash: fnv1a64:09e7583622d535b9
% canonical-lessons: FR-C394-benefice, FR-C394-perte, FR-C394-faillite, FR-C394-rembourser, FR-C394-emprunt

\chapter{Gagner et perdre}
\label{ch:fr-gagner-et-perdre}
\hlchaptermodality{394}

\begin{chapteropening}
\textbf{By the end of this chapter:} \emph{I can talk about business money: profit, a loss, going bankrupt, paying back and a loan.}

Complete the last lesson of chapter 394: I can talk about business money: profit, a loss, going bankrupt, paying back and a loan.
\end{chapteropening}

\section[le bénéfice]{le bénéfice — the profit; the benefit}

\label{lesson:FR-C394-benefice}



\begin{quote}
Before the new one: say the French for profitable, then the French for good value.
\end{quote}

\subsection*{You'll want to know: le bénéfice}
\textbf{le bénéfice} — \textquotedblleft{}the profit; the benefit\textquotedblright{}.

\begin{quote}
L'entreprise a fait un bénéfice cette année.
\end{quote}

\begin{quote}
\emph{The company made a profit this year.}
\end{quote}

\emph{Au bénéfice de} — for the benefit of.

\begin{cousinweb}
\emph{bénéfice} is a learned borrowing from Latin \emph{beneficium} \textquotedblleft{}a kindness, a favour\textquotedblright{} (EtymDB; Lewis \& Short), built from \emph{bene} \textquotedblleft{}well\textquotedblright{} and \emph{facere} \textquotedblleft{}to do\textquotedblright{}.
\end{cousinweb}

\begin{grammarlens}[title={What you've built}]
One more word for telling, explaining and describing at B1.
\end{grammarlens}

\subsection*{Your turn}
\begin{itemize}
\raggedright
  \item \emph{Say it:} \emph{le bénéfice}
  \item \emph{Say it:} \emph{le bénéfice}, once more
  \item \emph{Say it:} a sentence of your own with \emph{le bénéfice}
  \item \emph{Recall:} say the French for profitable, then the French for good value, then say \emph{le bénéfice} again
\end{itemize}

\begin{tcolorbox}[breakable,colback=teal!4,colframe=teal!35!black,title={Before you move on}]
What does le bénéfice mean? (\textquotedblleft{}The profit; the benefit\textquotedblright{}.) Say it once more.
\end{tcolorbox}
\section[la perte]{la perte — the loss}

\label{lesson:FR-C394-perte}



\begin{quote}
Before the new one: say the French for good value, then the French for the profit.
\end{quote}

\subsection*{You'll want to know: la perte}
\textbf{la perte} — \textquotedblleft{}the loss\textquotedblright{}.

\begin{quote}
Ce voyage est une perte de temps.
\end{quote}

\begin{quote}
\emph{This trip is a waste of time.}
\end{quote}

\emph{Une perte de temps} — a waste of time.

\begin{cousinweb}
\emph{perte} is inherited from Vulgar Latin \emph{perdita} \textquotedblleft{}a thing lost\textquotedblright{} (EtymDB), from \emph{perdere} \textquotedblleft{}to lose\textquotedblright{} (Lewis \& Short).
\end{cousinweb}

\begin{grammarlens}[title={What you've built}]
One more word for telling, explaining and describing at B1.
\end{grammarlens}

\subsection*{Your turn}
\begin{itemize}
\raggedright
  \item \emph{Say it:} \emph{la perte}
  \item \emph{Say it:} \emph{la perte}, once more
  \item \emph{Say it:} a sentence of your own with \emph{la perte}
  \item \emph{Recall:} say the French for good value, then the French for the profit, then say \emph{la perte} again
\end{itemize}

\begin{tcolorbox}[breakable,colback=teal!4,colframe=teal!35!black,title={Before you move on}]
What does la perte mean? (\textquotedblleft{}The loss\textquotedblright{}.) Say it once more.
\end{tcolorbox}
\section[la faillite]{la faillite — the bankruptcy}

\label{lesson:FR-C394-faillite}



\begin{quote}
Before the new one: say the French for the profit, then the French for the loss.
\end{quote}

\subsection*{You'll want to know: la faillite}
\textbf{la faillite} — \textquotedblleft{}the bankruptcy\textquotedblright{}.

\begin{quote}
Le restaurant a fait faillite après deux ans.
\end{quote}

\begin{quote}
\emph{The restaurant went bankrupt after two years.}
\end{quote}

\emph{Faire faillite} — to go bankrupt.

\begin{cousinweb}
\emph{faillite} was borrowed from Italian \emph{fallita} \textquotedblleft{}failure\textquotedblright{} (EtymDB), from Latin \emph{fallere} \textquotedblleft{}to deceive, to fail\textquotedblright{} (Lewis \& Short).
\end{cousinweb}

\begin{grammarlens}[title={What you've built}]
One more word for telling, explaining and describing at B1.
\end{grammarlens}

\subsection*{Your turn}
\begin{itemize}
\raggedright
  \item \emph{Say it:} \emph{la faillite}
  \item \emph{Say it:} \emph{la faillite}, once more
  \item \emph{Say it:} a sentence of your own with \emph{la faillite}
  \item \emph{Recall:} say the French for the profit, then the French for the loss, then say \emph{la faillite} again
\end{itemize}

\begin{tcolorbox}[breakable,colback=teal!4,colframe=teal!35!black,title={Before you move on}]
What does la faillite mean? (\textquotedblleft{}The bankruptcy\textquotedblright{}.) Say it once more.
\end{tcolorbox}
\section[rembourser]{rembourser — to pay back, to refund}

\label{lesson:FR-C394-rembourser}



\begin{quote}
Before the new one: say the French for the loss, then the French for the bankruptcy.
\end{quote}

\subsection*{You'll want to know: rembourser}
\textbf{rembourser} — \textquotedblleft{}to pay back, to refund\textquotedblright{}.

\begin{quote}
Je te rembourse demain, promis.
\end{quote}

\begin{quote}
\emph{I will pay you back tomorrow, I promise.}
\end{quote}

\emph{Se faire rembourser} — to get a refund.

\begin{cousinweb}
\emph{rembourser} is built from \emph{re-}, \emph{en-} and \emph{bourse} \textquotedblleft{}purse\textquotedblright{}: to put back into the purse. \emph{bourse} comes from Late Latin \emph{bursa} \textquotedblleft{}a hide, a leather bag\textquotedblright{} (Lewis \& Short).
\end{cousinweb}

\begin{grammarlens}[title={What you've built}]
One more word for telling, explaining and describing at B1.
\end{grammarlens}

\subsection*{Your turn}
\begin{itemize}
\raggedright
  \item \emph{Say it:} \emph{rembourser}
  \item \emph{Say it:} \emph{rembourser}, once more
  \item \emph{Say it:} a sentence of your own with \emph{rembourser}
  \item \emph{Recall:} say the French for the loss, then the French for the bankruptcy, then say \emph{rembourser} again
\end{itemize}

\begin{tcolorbox}[breakable,colback=teal!4,colframe=teal!35!black,title={Before you move on}]
What does rembourser mean? (\textquotedblleft{}To pay back, to refund\textquotedblright{}.) Say it once more.
\end{tcolorbox}
\section[l'emprunt]{l'emprunt — the loan (that you take out)}

\label{lesson:FR-C394-emprunt}



\begin{quote}
Before the new one: say the French for the bankruptcy, then the French for to pay back.
\end{quote}

\subsection*{You'll want to know: l'emprunt}
\textbf{l'emprunt} — \textquotedblleft{}the loan (that you take out)\textquotedblright{}.

\begin{quote}
Ils ont fait un emprunt pour acheter la maison.
\end{quote}

\begin{quote}
\emph{They took out a loan to buy the house.}
\end{quote}

From \emph{emprunter}, to borrow. Masculine.

\begin{cousinweb}
\emph{emprunt} is formed on \emph{emprunter} \textquotedblleft{}to borrow\textquotedblright{}, usually traced to a Vulgar Latin \emph{imprūmūtāre}, built on Latin \emph{mūtuum} \textquotedblleft{}a loan\textquotedblright{} (Lewis \& Short).
\end{cousinweb}

\begin{grammarlens}[title={What you've built}]
One more word for telling, explaining and describing at B1.
\end{grammarlens}

\subsection*{Your turn}
\begin{itemize}
\raggedright
  \item \emph{Say it:} \emph{l'emprunt}
  \item \emph{Say it:} \emph{l'emprunt}, once more
  \item \emph{Say it:} a sentence of your own with \emph{l'emprunt}
  \item \emph{Recall:} say the French for the bankruptcy, then the French for to pay back, then say \emph{l'emprunt} again
\end{itemize}

\begin{tcolorbox}[breakable,colback=teal!4,colframe=teal!35!black,title={Before you move on}]
What does l'emprunt mean? (\textquotedblleft{}The loan (that you take out)\textquotedblright{}.) Say it once more.
\end{tcolorbox}
